\documentclass{article}
\usepackage{amsmath, amssymb}
% Page geometry
\usepackage[margin=1in]{geometry}
\title{Orthogonality of OFDM Subcarriers}
\author{}
\date{}

\begin{document}

\maketitle
\section*{OFDM Subcarriers}
The $k^{th}$ subcarrier of an OFDM waveform can be expressed as:
\begin{equation}
    p_k(t) = p_{tx}(t) e^{j 2 \pi k F_{spacing} t},
\end{equation}
where $p_{tx}(t)$ is a rectangular pulse defined by:
\begin{equation}
    p_{tx}(t) = 
    \begin{cases} 
        1 & \text{for } 0 < t < T_{pulse}, \\
        0 & \text{otherwise}.
    \end{cases}
\end{equation}
Similarly, the $l^{th}$ subcarrier is given by:
\begin{equation}
    p_l(t) = p_{tx}(t) e^{j 2 \pi l F_{spacing} t}.
\end{equation}

\section*{Orthogonality Condition}
To prove that $p_k(t)$ and $p_l(t)$ are orthogonal for $F_{spacing} = \frac{1}{T_{pulse}}$ and $k \neq l$, we need to show that their inner product over one period equals zero.

The inner product is defined as:
\begin{equation}
    \int_0^{T_{pulse}} p_k(t) \cdot p_l^*(t) \, dt = \int_0^{T_{pulse}} p_{tx}(t) e^{j 2 \pi (k-l) F_{spacing} t} \, dt,
\end{equation}
where $p_{tx}^2(t) = p_{tx}(t)$ because it is a rectangular pulse with values of only 0 and 1.

Substituting $F_{spacing} = \frac{1}{T_{pulse}}$, the integral simplifies to:
\begin{equation}
    \int_0^{T_{pulse}} e^{j 2 \pi (k-l) \frac{1}{T_{pulse}} t} \, dt.
\end{equation}
For $k \neq l$, this integral evaluates to zero, indicating that $p_k(t)$ and $p_l(t)$ are orthogonal.

\end{document}
